% 由 scripts/publication/tables.py 自动生成，请勿手改。
% 需要宏包：booktabs（\usepackage{booktabs}）
\begin{table}[htbp]
  \centering
  \caption{Step 13 压力情景：交叉损耗放大 ×1/×2/×5/×10 的最大逐路损耗}
  \label{tab:step13-sensitivity}
  \begin{tabular}{lrrrr}
    \toprule
    方案 & ×1 (dB) & ×2 (dB) & ×5 (dB) & ×10 (dB) \\
    \midrule
    A & 6.3333 & 6.3333 & 6.3333 & 6.3333 \\
    R5U & 6.3222 & 6.3222 & 6.3222 & 6.3222 \\
    D0 & 5.3638 & 5.3638 & 5.3638 & 5.3638 \\
    D56 & 5.3462 & 5.3462 & 5.3462 & 5.3462 \\
    F5 & 5.9980 & 6.0117 & 6.0526 & 6.2857 \\
    F56 & 5.0783 & 5.1038 & 5.1803 & 5.9941 \\
    \midrule
    A & 6.5773 & 6.5773 & 6.5773 & 6.5773 \\
    R5U & 6.5780 & 6.5780 & 6.5780 & 6.5780 \\
    D0 & 5.6207 & 5.6207 & 5.6207 & 5.6207 \\
    D56 & 5.6217 & 5.6217 & 5.6217 & 5.6217 \\
    F5 & 6.4459 & 6.4970 & 6.6503 & 8.6653 \\
    F56 & 5.4995 & 5.5452 & 5.7264 & 8.4333 \\
    \bottomrule
  \end{tabular}
  \par\vspace{2pt}
  \begin{flushleft}\footnotesize
  固定几何复算：仅把 <20° 区间交叉损耗放大，再重算全部逐路损耗；不重新布线，路线选择不变。 \\
  ×5/×10 是实验压力情景，不代表已验证的真实物理误差范围。 \\
  \end{flushleft}
\end{table}
